\title{LongRoPE: Extending LLM Context Window Beyond 2 Million Tokens}

\begin{abstract}
An extensive context window represents a valuable capability for large language models (LLMs). Nonetheless, contemporary approaches to enlarging context windows typically reach up to approximately 128k tokens, constrained by the high expense associated with fine-tuning, the limited availability of lengthy textual datasets, and problematic token values that arise from unfamiliar positional encodings.
\end{abstract}

This study presents LongRoPE, which, for the first time, expands the context window of pre-trained large language models (LLMs) to an impressive \textbf{2048k} tokens, requiring as few as 1k fine-tuning steps on sequences up to 256k tokens, while preserving the model’s capabilities on its initial, shorter context window. This advancement is the result of three core contributions: (i) we systematically uncover and leverage two specific non-uniformities in positional interpolation using a rapid search method, leading to improved fine-tuning initialization and permitting up to an 8$\times$ increase in context length even without fine-tuning; (ii) we propose a stepwise extension approach, whereby a LLM is first fine-tuned to 256k tokens, followed by a secondary positional interpolation on the newly extended model, reaching a 2048k context window; (iii) we recalibrate LongRoPE at 8k sequence length to restore the model’s original short-context performance. Comprehensive experiments with LLaMA2 and Mistral on diverse benchmarks showcase the robust effectiveness of our approach. Models extended using LongRoPE retain the foundational model architecture, with only slight adjustments to positional embeddings, allowing the majority of existing optimizations to remain applicable. The codebase will be made publicly accessible at \texttt{https://github.com/microsoft/LongRoPE}

\section{Introduction}

Although Large Language Models (LLMs) have achieved impressive results across a range of tasks~\cite{gpt4,llama2}, they are often constrained by their finite context window sizes; for example, LLaMA2 is restricted to processing up to 4096 tokens~\cite{llama2}. When inputs exceed this window, the performance of LLMs deteriorates, largely because the model is not accustomed to handling positions outside of those encountered during training. Such limitations become particularly problematic in applications such as in-context learning with large numbers of examples~\cite{Huang2023FewerIM}, as well as in the operation of LLM agents~\cite{park2023generative,madaan2023selfrefine}.

Recent studies have demonstrated that it is possible to extend the context window of a pre-trained LLM to approximately 128k tokens through fine-tuning on longer sequences~\cite{longlora,pi,yarn,zhang2024soaring,liu2023scaling}. However, several significant barriers hinder further expansion of the context window. Firstly, assigning previously unused position indices during extension often results in numerous aberrant values, causing out-of-distribution complications and impeding convergence during fine-tuning~\cite{pi}. This issue is exacerbated when increasing the window from 4k to over 1000k positions, where more than 90\% of the positions are newly introduced. Secondly, fine-tuning typically demands the availability of texts matching the extended length, but datasets with texts longer than 1000k tokens are scarce. Additionally, training with such extremely long documents requires immense computational resources, making the process both time- and resource-intensive. Thirdly, expanding the context window to extreme lengths tends to dilute attention over many token positions, which compromises the model’s ability to perform well on tasks involving shorter contexts~\cite{pi}.

To address the first challenge, one method is to apply interpolation to RoPE positional embeddings~\cite{rope,pi}, which effectively rescales newly encountered position indices so that they fit within the pretrained positional range, as depicted in Fig.~\ref{fig:overview}. The Position Interpolation (PI) technique~\cite{pi} achieves this by performing a linear interpolation of the rotary angles in RoPE based on the degree of context extension. In contrast, NTK~\cite{ntk,dynamicntk} introduces non-uniform strategies, involving both unequal interpolation and extrapolation along different RoPE dimensions. Meanwhile, YaRN~\cite{yarn} further refines this idea by dividing RoPE dimensions according to their frequencies, implementing extrapolation, NTK scaling, or linear interpolation within each respective group. Nonetheless, within the Transformer model, the positional embedding inherently possesses a \emph{complex and heterogeneous information entropy} distribution. Current methods fail to sufficiently exploit these intricate non-uniformities, resulting in loss of information and thus constraining the achievable context window.

Section~\ref{sec:analysis} presents two primary empirical observations: \textbf{(1)} Successful positional interpolation must account for two distinct types of non-uniformity, specifically those arising from different RoPE dimensions and the positions of tokens. Reduced interpolation is more suitable for dimensions with lower RoPE indices and for tokens at earlier positions; however, the most effective strategies vary according to the desired extended sequence length. \textbf{(2)} Incorporating these non-uniformities within the positional interpolation process allows for better preservation of the original RoPE information, with particular emphasis on essential dimensions and initial token locations. This approach curtails information loss from interpolation, thereby yielding a superior starting point for subsequent fine-tuning. Additionally, it makes it possible to extend sequences up to 8$\times$ their original length even without further fine-tuning.

Building upon these observations, we present \textbf{LongRoPE}, a robust approach designed to expand the LLM context window to exceed 2 \emph{million} tokens. The LongRoPE framework is founded on three central advancements. Firstly, {LongRoPE} leverages the multidimensional and non-uniform nature of positional interpolation to its full potential. Specifically, it determines optimal rescale factors for the rotation angles in RoPE across all dimensions by taking token positions into account. Given that the parameter space for identifying these rescale factors grows exponentially with larger extension ratios, {LongRoPE} employs an evolutionary search strategy incorporating two optimization methods that significantly improve the efficiency of the search process. Figure~\ref{fig:overview} illustrates a sample result of the optimized and rescaled RoPE.

Subsequently, {LongRoPE} employs a resource-efficient, incremental expansion method that enables a 2048k context window, circumventing the necessity to fine-tune directly on extremely lengthy texts, which are both rare and difficult to source. The approach starts with identifying a 256k context length on the pre-trained LLM and conducting fine-tuning at this scale. Given that our approach to non-uniform positional interpolation permits up to an 8$\times$ increase in the context window without further fine-tuning, a second search is then performed for optimal RoPE rescale factors using the previously fine-tuned, expanded LLM. Through this process, the 2048k context window is successfully realized for LLaMA2 and Mistral~\cite{mistral}.

In order to address potential drops in performance with the original, shorter context window, {LongRoPE} maintains optimization of the RoPE rescale factors on the expanded LLM. In a manner analogous to increasing the context length from 256k to 2048k, our search algorithm is utilized to decrease the context window to 4k and 8k on the LLM fine-tuned at 256k, thus promoting minimal positional interpolation. At inference time, when the input sequence length falls below 8k, RoPE is modified using the identified optimal rescale factors.

A wide range of experiments conducted on multiple LLMs and diverse long-context tasks validate the efficacy of our proposed approach. Our results indicate that LongRoPE consistently preserves low perplexity across evaluation lengths spanning from 4k up to 2048k, surpasses 90\% accuracy in passkey retrieval, and matches the performance of baselines on typical benchmarks situated within a 4096-token context window. Furthermore, LongRoPE is broadly compatible with any LLMs utilizing RoPE embedding. We intend to make both our codebase and the LongRoPE-2048k models publicly available.

\section{Non-uniformity in Positional Interpolation}
\label{sec:analysis}

\subsection{Preliminary}

Transformer architectures depend on explicit positional cues, typically supplied through position embeddings, to capture the sequence of input tokens. In this study, we concentrate on the RoPE~\cite{rope} positional embedding mechanism, a technique extensively adopted in state-of-the-art LLMs. For a token located at index $n$, its associated RoPE encoding can be succinctly expressed as:

\begin{equation}
		\fontsize{7.8}{7.8} \selectfont
	\label{eq:rope}
	\left[ cos(n\theta_0), sin(n\theta_0), cos(n\theta_1), \cdot\cdot\cdot, cos(n\theta_{d/2-1}), sin(n\theta_{d/2-1}) \right]   
\end{equation}

where $d$ denotes the embedding dimension, $n\theta_i$ corresponds to the rotary phase applied to the token located at position $n$, and $\theta_i=\theta^{-2i/d}$ encodes the rotational frequencies. In the case of RoPE, $\theta$ is typically set to a base value of 10000.

\textbf{\emph{Context window extension ratio $s$} and positional interpolation}. The extension ratio $s$ is defined as the quotient of the new (extended) context length $L'$ by the original context length $L$, that is, $s=\frac{L'}{L}$.

To increase the context window from $L$ to $L'$, existing positional interpolation techniques typically involve reducing the rotation frequencies $\theta_i$ according to the expansion factor $s$. Defining $\beta=\theta^{2/d}$ and using $\lambda$ to represent the genuine scaling parameter linked to $s$, we present a unified framework encompassing these positional interpolation strategies:

\begin{equation}	
	\fontsize{7.5}{7.5} \selectfont
	\label{eq:unified_rope}
	\begin{aligned}
		\left[cos\left(\frac{n}{\lambda(\beta)^0}\right), sin \left(\frac{n}{\lambda(\beta)^0}\right), cos\left(\frac{n}{\lambda(\beta)^1}\right), \cdot\cdot\cdot,  sin \left(\frac{n}{\lambda(\beta)^{d/2-1}}\right) \right]   
	\end{aligned}
\end{equation}

\textbf{Linear positional interpolation (PI)}. The PI approach~\cite{pi} proposes interpolating position indices linearly for inputs within the model's original sequence length. When extending the sequence length by a factor of $s$, the rotation angles corresponding to every position are uniformly scaled down by $\lambda=s$ for all RoPE dimensions. This linear compression causes position representations to cluster more closely together, making it more challenging for the model to discern tokens that are near each other in the sequence. As a result, PI's effectiveness generally declines as the extension ratio increases.

\textbf{NTK-based interpolation and extrapolation}. ~\cite{ntk,dynamicntk} investigate RoPE from the standpoint of information representation and leverage Neural Tangent Kernel (NTK) theory~\cite{jacot2018neural,tancik2020fourier} for analysis. To address the problem of positional crowding inherent in PI, they propose reallocating the interpolation burden among the different RoPE dimensions. This is achieved by reducing the scaling of lower (high frequency) dimensions while amplifying the scaling for higher (low frequency) ones, thereby facilitating both interpolation and extrapolation across positions, where $\lambda=s^{i}$. The enhanced dynamic NTK~\cite{dynamicntk} further adapts the extension factor at each position depending on the sequence length in use. In contrast to PI, which relies on additional fine-tuning, NTK-informed approaches are capable of lengthening context windows without requiring fine-tuning, although such extension is typically limited to a maximum factor of 4$\times$.

\textbf{YaRN}~\cite{yarn} presents a notable advancement in the efficacy of positional interpolation. The method categorizes RoPE dimensions into three distinct groups according to their associated frequencies, and applies unique interpolation techniques to each. Specifically, extrapolation with $\lambda$=1 is utilized for dimensions with high frequencies, linear interpolation (PI) is reserved for those with low frequencies, while dimensions with intermediate frequencies are processed using the NTK method. A central feature of YaRN is this frequency-based partitioning of RoPE dimensions, which is currently determined through manual, empirical analysis. As a result, this approach may lead to less-than-optimal outcomes when adapting to novel LLMs.

\subsection{Study on Non-uniform Positional Interpolation}

Drawing from insights in NTK and YaRN, it is evident that their performance improvements stem from incorporating non-linear mechanisms—most notably by treating distinct frequencies within RoPE dimensions with tailored strategies for interpolation and extrapolation. Nonetheless, existing approaches to non-linearity largely depend on manually crafted heuristics. This observation prompts two central inquiries: (1) Does the current approach to positional interpolation represent the optimal solution? (2) Could there be other, as yet undiscovered, forms of non-linearity?

To address these questions, we employ evolution search (refer to Sec.~\ref{sec:method}) to identify improved non-uniform positional interpolations for LLaMA2-7B. This process is steered by perplexity measurements, evaluated over 5 randomly selected samples from the PG19~\cite{pg19} validation dataset. Our experimental investigation yields several important insights.

\textbf{Finding 1}: \textit{RoPE dimensions exhibit substantial non-uniformities, which are not effectively handled by current positional interpolation methods.}

We optimize the value of $\lambda$ for each RoPE dimension as specified in Eq.~\ref{eq:unified_rope}. Table~\ref{tbl:exp1} presents the perplexity results of LLaMA2-7B, evaluated on the PG19 and Proof-pile~\cite{proof-pile} test sets using different approaches, all without additional fine-tuning. Our optimized configuration delivers substantial perplexity reductions, indicating that both the existing linear (PI) and non-uniform (Dynamic-NTK, YaRN) interpolation strategies do not yield the best results. It is particularly noteworthy that YaRN lags behind both PI and NTK on the PG19 benchmark, likely due to its failure to maintain the intended context window length in settings without further fine-tuning. As an illustration, when the context size reaches 8k, YaRN’s perplexity sharply increases beyond the 7k mark.

In our investigation, we observed that the rescaled coefficients $\lambda$ in Eq.~\ref{eq:unified_rope} are non-uniform, contrasting with the constant scaling $s$ used in PI, the formula in NTK, and YaRN’s approach of group-wise rescaling. The introduction of these non-uniform scaling factors leads to notable improvements in LLaMA2’s language modeling capabilities (measured by perplexity) for 8k and 16k context lengths, and this enhancement occurs without the need for additional fine-tuning. The primary reason for this improvement is that the induced positional embeddings closely retain the properties of the original RoPE—most importantly, those of key dimensions—thus alleviating the challenge LLMs face in differentiating nearby token positions.

\textbf{Finding 2}: \textit{RoPE for the initial tokens in the input sequence should be extrapolated with less interpolation.} 

We propose that for the first $\hat{n}$ tokens in the input sequences, their RoPE should be subjected to minimal interpolation. This is based on the observation that these initial tokens typically acquire higher attention weights, making them particularly influential within attention layers—a finding consistent with previous work in Streaming LLM~\cite{streamingllm} and LM-Infinite~\cite{han2023lminfinite}. To test this hypothesis, we enlarge the context window to 8k and 16k via PI and NTK, while varying the number of leading tokens ($\hat n$ = 0, 2, ..., 256) for which position interpolation is not applied. Notably, setting $\hat n = 0$ yields the baseline PI and NTK models. Table~\ref{tbl:tokenposition} elucidates two main findings: \textbf{(1)} exempting the initial tokens from position interpolation consistently enhances model performance, and \textbf{(2)} the choice of optimal $\hat n$ is contingent upon the specific target window size.

\textbf{Finding 3}: \textit{Non-uniform positional interpolation effectively extends LLM context window in both fine-tuning and non-fine-tuning settings.} 

Although our investigated non-uniform position interpolation markedly enhances extension performance at 8k and 16k without any additional training, extending to longer contexts necessitates fine-tuning. Consequently, we perform fine-tuning on LLaMA2-7B with our identified RoPE configuration for a 64k context window (refer to Appendix for experimental details). According to Table~\ref{tbl:finetune}, our approach consistently surpasses PI and YaRN, whether or not fine-tuning has been applied to LLaMA2-7B. This performance gain can be attributed to the effective implementation of non-uniform positional interpolation, which reduces information loss and yields superior initialization for the fine-tuning process.

\textbf{Summary}. This work identifies two types of non-uniformities—variation in RoPE dimensions and disparities in token positions. Strategic exploitation of these factors in positional interpolation demonstrably enhances context extension capabilities for LLMs.

\section{LongRoPE}

\label{sec:method}
Building on these insights, we propose LongRoPE—a method that initially employs a novel, efficient search algorithm designed to harness both non-uniformities. Leveraging this approach, LongRoPE is able to expand the LLM context window to accommodate over 2 million tokens.

\subsection{Problem Formulation}

The presence of these two forms of non-uniformity significantly expands the solution landscape and adds layers of difficulty to the optimization process. To tackle this challenge, we reformulate the multidimensional non-uniform position interpolation optimization as a search problem.

Consider an LLM designed for a context window size of $L'$ processing long input documents $\mathbf{X}$, where every sequence $\mathbf{x} \in \mathbf{X}$ contains more tokens than $L'$. We express the initial rotary angle for the $i^{th}$ dimension in the RoPE embedding at token position $n$ as $\frac{n}{\beta^i}$. Based on this setup, the optimization problem can be described as:

\begin{equation}
\begin{gathered}
		\label{eq:problem}

\underset {\mathbf{x}\in\mathbf{X};\, |\mathbf{x}|\ge L'}{ \operatorname {arg\,min} } \,  \mathcal{L}\, \big( \text{LLM} (\text{RoPE}, \mathbf{X}) \big),\quad
\\
\scalebox{0.93}{
$\underset {\begin{subarray}{l} i=0,\ldots,\frac{d}{2}-1;\\ n\in[0,|\mathbf{x}|);\end{subarray}}{\text{RoPE}_{(n)}} = \left[\,\cdots,\cos\left(\mathbb{I}(\hat{\lambda}_{i}, \hat{n})\frac{n}{\beta^i}\right),\ \sin\left(\mathbb{I}(\hat{\lambda}_{i}, \hat{n})\frac{n}{\beta^i}\right),\cdots\right] $}
\\
\text{where } \mathbb{I}(\hat{\lambda}_{i},\hat{n})=
\begin{cases}
1 & \text{if } n < \hat{n}\\
\frac{1}{\lambda_i} & \text{if } n \geq \hat{n}
\end{cases}
\end{gathered}
\end{equation}

In this approach, we introduce a collection of scaling factors, $\mathbb{I}(\hat{\lambda}_{i},\hat{n})$, aimed at addressing two distinct sources of non-uniformity. Here, $\hat{\lambda_i}$ represents variation across RoPE dimensions, while $\hat{n}$ pertains to positional irregularities among tokens. Specifically, the term $\mathbb{I}(\hat{\lambda}_{i},\hat{n})$ is employed to adjust the rotation angle corresponding to the $i^{th}$ RoPE dimension, using $\hat\lambda_{i}$ as the scaling coefficient and $\hat{n}$ as a positional cutoff. For the first $\hat{n}$-1 tokens, the adjustment factor $\hat\lambda_{i}$ is not applied, retaining the standard RoPE rotary angle $\frac{n}{\beta^i}$. When the token position satisfies $n\ge\hat{n}$, the rescaling factor becomes active.

For a specified target context window length $L'$, the task is to determine the optimal set of rescale factors ($\mathbb{I}(\hat{\lambda}_{0},\hat{n})$, $\mathbb{I}(\hat{\lambda}_{1},\hat{n})$, ... $\mathbb{I}(\hat{\lambda}_{i},\hat{n})$, ...) across the $1^{st}$ through $d^{th}$ dimensions of RoPE. The aim is that, upon applying these rescale factors to the target $\text{LLM}$'s $\text{RoPE}$, the model will yield the lowest possible next-token prediction loss, $\mathcal{L}$ (measured by perplexity), when provided with input sequences $\mathbf{X}$ of length $L'$.

\subsection{Searching the Non-uniform Position Interpolation}

To address the problem formulated in Eq.~\ref{eq:problem}, we propose an efficient and straightforward approach that identifies optimal RoPE rescale factors, thereby taking full advantage of the multidimensional irregularities present in position embedding.

\textbf{Search space}. We construct a comprehensive search space that accommodates both forms of non-uniformity. An overview of this search space is provided in Table~\ref{tbl:searchspace}. More precisely, we permit the assignment of an individualized rescale factor to each RoPE embedding dimension. To facilitate a simpler search space, we search over the variables $\lambda_i$ and $\hat{n}$ instead of $\mathbb{I}(\hat{\lambda}_{i},\hat{n})$, where  $\hat{\lambda}_{i}=1/\lambda_i$. 
As depicted in Table~\ref{tbl:searchspace}, the parameter $\lambda_i$ is explored in the interval from a minimum of 1.0 (corresponding to direct extrapolation) up to a maximum of $s\times1.25$ (which extends interpolation beyond PI), with increments of 0.01. Here, $s$ denotes the desired ratio for extending the context window.

The parameter $\hat{n}$ specifies how many leading token positions utilize the unaltered RoPE embedding, thereby bypassing position interpolation. In practice, we let $\hat{n}$ vary among the values \{0, 1, 2, 4, 8, 12, 16, 20, 24, 28, 32, 64, 128, 256\}. If $\hat{n}=0$, then every token position employs the optimized rescale factors.

\textbf{Evolution-based search}. The search space detailed in Table~\ref{tbl:searchspace} encompasses a vast array of positional interpolation strategies, making comprehensive exploration computationally prohibitive. For instance, considering an extension factor of $s=4\times$ yields $400^{128/2}\times14=4\times10^{167}$ possible configurations. As the extension ratio increases, the size of the search space grows exponentially. To make the exploration tractable, we adopt evolution search~\cite{guo2020single} and incorporate two optimization strategies that substantially enhance the efficiency of the search process. The overall methodology is outlined in Algorithm~\ref{alg:search}.

\textit{Optimized initial population generation}. Rather than relying solely on random initialization for the population of $P$ rescale factors, we incorporate the three RoPE rescale factors associated with PI, NTK, and YaRN directly as distinct individuals in the starting population. The rest of the individuals, numbering $P$-3, are produced by applying random mutations to these three base rescale factors with a probability of $p$.

\textit{Monotonically non-decreasing constraint}. Once the initial population has been generated, we evaluate the perplexity of the LLM for every individual. This involves applying the associated RoPE rescale factors to the target LLM and determining the perplexity for the input $\mathbf{X}$. The k individuals with the highest performance are chosen as parents for the subsequent evolutionary steps. Nonetheless, because the search space is enormous, simple mutation and crossover operations may result in the evaluation of low-quality solutions, incurring excessive computations of perplexity. This issue becomes even more pronounced for large $L'$, as each perplexity assessment is computationally expensive.

To tackle this issue, we enforce a monotonic non-decreasing restriction on the sequence of sampled RoPE rescaling coefficients such that $\lambda_i \le \lambda_{i+1}$. Only those RoPE configurations that satisfy this monotonicity constraint are employed during perplexity assessment in LLMs, thereby greatly curtailing the computational overhead associated with the search process. In particular, we stipulate that $\lambda_i$ progresses in a monotonically increasing manner across RoPE dimensions, with $i$ ranging from $0$ to $63$. This requirement for monotonicity across dimensions is inspired by insights from NTK theory~\cite{jacot2018neural,tancik2020fourier,ntk}, which indicates that dimensions with higher frequencies (corresponding to lower indices) should undergo less interpolation (meaning a smaller value of $\lambda_i$), while those with lower frequencies (higher indices) can accommodate greater interpolation (therefore, a larger $\lambda_i$).

\begin{algorithm}[t]
	\small
	\caption{The search algorithm}
	\textbf{Input:} target LLM, input samples $\mathbf{X}$, population size $P$, mutation size $N_1$, crossover size $N_2$, maximum number of iterations $\mathcal{T}$, mutation probability $p$\\

	\begin{algorithmic}[1]
		\label{alg:search}
		\STATE Top-k $\leftarrow \phi$;
		\STATE $\text{P}_0 \leftarrow$ \textit{Initialize\_population\_with\_optimization} ($P$, $\mathbf{X}$, $p$);
		\FOR{$i=1$ to $\mathcal{T}$}
		\STATE \textit{Compute\_perplexity} (LLM, $\text{P}_{i-1}$, $\mathbf{X}$);
		\STATE Top-k $\leftarrow$ \textit{Update\_Topk} (Top-k, $\text{P}_{i-1}$);
		\STATE $\text{P}_{mutation} \leftarrow$ \textit{Mutation\_with\_mono\_constraint} (Top-k, $N_1$, $p$);
		\STATE $\text{P}_{crossover} \leftarrow$ \textit{Crossover\_with\_mono\_constraint} (Top-k, $N_2$);
		\STATE $\text{P}_i \leftarrow \text{P}_{mutation} \cup \text{P}_{crossover} \cup$ Top-k;
		\ENDFOR
		\STATE Return the individual within Top-k that has the minimum perplexity;
	\end{algorithmic}
\end{algorithm}

\textbf{8$\times$ extension without fine-tuning}. Through the use of evolutionary search, our approach successfully discovers optimal non-uniform RoPE rescaling factors, maintaining essential dimensions and position encodings to reduce information degradation from interpolation. As shown in Fig.\ref{fig:non-ft}, our strategy enables LLaMA2 to expand its context window from 4k up to 32k tokens without any need for fine-tuning. By comparison, other techniques like PI as well as non-uniform NTK and YaRN exhibit substantial increases in perplexity once the context length surpasses a 2$\times$ extension.

\subsection{Extending LLM Context Window to 2048K}
\label{sec:extendto2048k}

\textbf{Progressive extension to 2048k}. In this section, we present our approach for expanding the context window of pre-trained LLMs from the typical 4k tokens up to beyond 2048k tokens. As our results indicate, utilizing non-uniform positional interpolation enables an extension up to 8$\times$ the original window without the need for fine-tuning. However, achieving a significantly larger extension—on the order of 512$\times$—necessitates fine-tuning. One possible approach involves searching for optimal RoPE rescale factors at the desired 2048k context length, followed by fine-tuning. This, however, is impeded by the extremely high computational cost associated with such training. Furthermore, our empirical observations suggest that effective fine-tuning of LLMs becomes increasingly difficult as the extension ratio grows large (refer to Appendix).

Luckily, LongRoPE proves to be robust across both base and fine-tuned extended LLMs. In light of this, we propose a streamlined and incremental strategy that reaches the desired 2048k sequence length using only 1,000 fine-tuning iterations, all conducted with a maximum training length of 256k.

$\Diamond$ \textit{LongRoPE search for adapting pre-trained LLM to 256k}. Using LLaMA2 as a case study, we perform a search aiming for context window lengths of 128k and 256k. This corresponds to expansion factors of 32$\times$ and 64$\times$ at this phase.

$\Diamond$ \textit{Fine-tuning to 256k}. In this stage, we adapt the pre-trained LLM to support a context window length of 256k. The process begins by fine-tuning LLaMA2 for 400 steps utilizing RoPE rescaled factors set for 128k. After completing this phase, we switch the RoPE rescaled factors to correspond to 256k on the resulting checkpoint and continue with 600 more steps of fine-tuning. This two-stage approach is more efficient than beginning fine-tuning directly at 256k.

$\Diamond$ \textit{Scaling fine-tuned extended LLM to 2048k using LongRoPE search}. In the final stage, we conduct an additional search step on the LLM previously fine-tuned to a 256k context length. Through this process, we successfully achieve an exceptionally large context window of 2048k, accomplished without any further fine-tuning. This yields a total extension factor of $512\times$.

\textbf{Recovery in shorter context windows}. When scaling up to an exceptionally large 2048k context window, we observe degraded performance within the initial context length. This phenomenon is well-documented with positional interpolation~\cite{pi}, as it compels the position embeddings associated with the original context window to cluster densely in a limited area of the embedding space, which diminishes the model's effectiveness. Specifically, applying a 512$\times$ extension factor causes the embeddings for positions in the standard 4k context window to be tightly packed.

To address this, an additional evolution search is conducted on the expanded LLM, specifically to fine-tune the RoPE rescale factors for shorter context lengths (such as 4k and 8k tokens). The upper limit for the searched $\lambda$ is lowered, reflecting the reduced need for positional interpolation at these shorter lengths. At inference time, the LLM automatically adapts the RoPE rescale factors based on the context length in use.

\section{LongRoPE}

\label{sec:method}
Building on these observations, we propose LongRoPE. This method initially incorporates an effective search procedure designed to leverage both forms of non-uniformity, and subsequently applies this approach to expand the context window of LLMs to over 2 million tokens.

\subsection{Problem Formulation}

These two sources of non-uniformity expand the solution landscape considerably and complicate the optimization process. To manage this complexity, we recast the multidimensional non-uniform position interpolation optimization challenge into a search-based formulation.

Considering a large language model (LLM) with a context window of size $L'$ processing long-form input documents $\mathbf{X}$, where each sequence $\mathbf{x} \in \mathbf{X}$ contains more than $L'$ tokens, let us represent the initial rotary angle for the $i^{th}$ dimension at position $n$ in RoPE embeddings as $\frac{n}{\beta^i}$. The corresponding optimization task is structured as follows:

\begin{equation}
\begin{gathered}
		\label{eq:problem}

\underset {\mathbf{x}\in\mathbf{X}; \, |\mathbf{x}|\ge L'}{ \operatorname {arg\,min} } \,  \mathcal{L}\, \left( \text{LLM} (\text{RoPE}, \mathbf{X})\right),\quad \text{such that}
\\
\scalebox{0.93}{
$\underset {\begin{subarray}{l} i=0,\cdot\cdot,\frac{d}{2}-1;\\ n\in[ 0,|\mathbf{x}|);\end{subarray}}{\text{RoPE}_(n)}= {\left[\cdots, \cos\left(\mathbb{I}(\hat{\lambda}_{i}, \hat{n})\cdot\frac{n}{\beta^i}\right),\ \sin\left(\mathbb{I}(\hat{\lambda}_{i}, \hat{n})\cdot\frac{n}{\beta^i}\right),\cdots\right] }$}\\
	\text{with} \quad \mathbb{I}(\hat{\lambda}_{i},\hat{n})=\begin{cases}
	1&\text{if $n<\hat{n}$}\\
{\frac{1}{\lambda}_{i}}&\text{if $n\geq\hat{n}$}
\end{cases}
	\end{gathered}
\end{equation}

In this approach, we define a collection of scaling factors, $\mathbb{I}(\hat{\lambda}_{i},\hat{n})$, designed to account for the two types of non-uniformities present. Here, $\hat{\lambda}_i$ corresponds to non-uniformity in RoPE dimensions, while $\hat{n}$ captures positional non-uniformity among tokens. The term $\mathbb{I}(\hat{\lambda}_{i},\hat{n})$ is specifically utilized to modify the rotation angle associated with the $i^{th}$ dimension of RoPE; $\hat\lambda_i$ serves as the adjustment coefficient, and $\hat{n}$ establishes the threshold for token positions. For the first $\hat{n}$ token positions, the scaling factor $\hat\lambda_i$ remains inactive, meaning the conventional RoPE rotary angle $\frac{n}{\beta^i}$ continues to be used. Once the position reaches or exceeds $n\ge\hat{n}$, the scaling factor is incorporated into the calculation.

Assuming a desired context window of size $L'$, the goal is to determine the most suitable set of rescale factors ($\mathbb{I}(\hat{\lambda}_{0},\hat{n})$, $\mathbb{I}(\hat{\lambda}_{1},\hat{n})$, ..., $\mathbb{I}(\hat{\lambda}_{i},\hat{n})$, ...) across the $1^{st}$ through $d^{th}$ RoPE dimensions. By applying these rescaled factors to the target $\text{LLM}$'s $\text{RoPE}$, we seek to minimize the next-token prediction loss, denoted as $\mathcal{L}$ (representing perplexity), on input data $\mathbf{X}$ containing $L'$ tokens.

\subsection{Searching the Non-uniform Position Interpolation}

To address the challenge outlined in Eq.~\ref{eq:problem}, we propose an approach that efficiently identifies optimal RoPE rescale factors, enabling the model to capitalize on the complex, multidimensional irregularities within position embeddings.

\textbf{Search space}. We construct an expansive search space that captures both forms of non-uniformity, as summarized in Table~\ref{tbl:searchspace}. In our formulation, each RoPE embedding dimension is assigned its own learnable rescale factor. To streamline the design, the search operates over the parameters $\lambda_i$ and $\hat{n}$, rather than directly over $\mathbb{I}(\hat{\lambda}_{i},\hat{n})$, where $\hat{\lambda}_{i}=1/\lambda_i$. According to Table~\ref{tbl:searchspace}, the range for $\lambda_i$ extends from a lower bound of 1.0 (which represents standard extrapolation) up to $s\times1.25$ (corresponding to interpolation regions beyond what PI offers), incremented in steps of 0.01; here, $s$ stands for the intended scaling ratio of the context window.

The parameter $\hat{n}$ specifies how many of the initial token positions retain the original RoPE embedding, foregoing position interpolation. In our experiments, we permit $\hat{n}$ to take values from the set \{0, 1, 2, 4, 8, 12, 16, 20, 24, 28, 32, 64, 128, 256\}. If $\hat{n}=0$, this means that all token positions employ the optimized rescale factors.

\textbf{Evolution-based search}. The range of options outlined in Table~\ref{tbl:searchspace} for positional interpolation encompasses a vast and complex search space, making efficient navigation a formidable task. As an illustration, performing an extension with $s=4\times$ results in $400^{128/2}\times14=4\times10^{167}$ possible configurations. Increasing the extension ratio further escalates the search space at an exponential rate. To manage this complexity, we employ evolution search~\cite{guo2020single} and integrate two additional optimization strategies to substantially enhance the efficiency of the search process. The complete procedure is described in Algorithm~\ref{alg:search}.

\textit{Optimized initial population generation}. Rather than generating all $P$ rescale factors at random for the initial population, we explicitly include the three RoPE rescale factors associated with PI, NTK, and YaRN among the starting individuals. The other $P$-3 members of the population are created by applying random mutations to these three rescale factors, each with a mutation probability $p$.

\textit{Monotonically non-decreasing constraint}. Once the initial population is generated, we evaluate each individual by calculating their LLM perplexity. To do this, the specified RoPE rescale factors are applied to the target LLM, and perplexity is measured for the input $\mathbf{X}$. The highest-performing $k$ individuals are then chosen as parents for the evolutionary process. Nonetheless, due to the enormous size of the search space, indiscriminate application of mutation and crossover operations may lead to the selection of suboptimal candidates, resulting in an excessive number of perplexity evaluations. This challenge is exacerbated for larger values of $L'$, as each perplexity computation requires significant inference time.

To tackle this issue, we enforce a monotonic non-decreasing constraint on the set of rescaled RoPE factors so that $\lambda_i \leq \lambda_{i+1}$ for all indices $i$. We only consider RoPE configurations that meet this monotonicity criterion when evaluating LLM perplexity, thus greatly diminishing the computational overhead of the search process. In particular, we stipulate that the sequence $\lambda_i$ must be monotonically increasing with respect to the RoPE dimension, i.e., for $i = 0, \ldots, 63$. This requirement is motivated by Neural Tangent Kernel (NTK) theory~\cite{jacot2018neural,tancik2020fourier,ntk}, which implies that dimensions associated with higher frequency content (smaller $i$) benefit from less interpolation (smaller values of $\lambda_i$), whereas dimensions with lower frequencies (larger $i$) can accommodate greater interpolation (larger $\lambda_i$).

\begin{algorithm}[t]
	\small
	\caption{The search algorithm}
	\textbf{Input:} target LLM, input samples $\mathbf{X}$, population size $P$, mutation size $N_1$, crossover size $N_2$, maximum iterations $\mathcal{T}$, mutation probability $p$\\

	\begin{algorithmic}[1]
		\label{alg:search}
		\STATE Initialize Top-k as $\phi$;
		\STATE Set initial population $\text{P}_0$ using \textit{Initialize\_population\_with\_optimization} ($P$, $\mathbf{X}$, $p$);
		\FOR{$i$ from $1$ to $\mathcal{T}$}
		\STATE Execute \textit{Compute\_perplexity} (LLM, $\text{P}_{i-1}$, $\mathbf{X}$);
		\STATE  Refresh Top-k with \textit{Update\_Topk} (Top-k, $\text{P}_{i-1}$);
		\STATE Generate mutated population $\text{P}_{mutation}$ using \textit{Mutation\_with\_mono\_constraint} (Top-k, $N_1$, $p$); 
		\STATE Generate crossover population $\text{P}_{crossover}$ via \textit{Crossover\_with\_mono\_constraint} (Top-k, $N_2$);
		\STATE Form next generation $\text{P}_i$ as $\text{P}_{mutation}$ $\cup$ $\text{P}_{crossover}$ $\cup$ Top-k;
		\ENDFOR
		\STATE Output the candidate with the minimum perplexity from Top-k;
	\end{algorithmic}
\end{algorithm}

\textbf{8$\times$ extension without fine-tuning}. By leveraging an evolutionary search approach, we successfully discover optimal non-uniform RoPE rescaling factors, selectively retaining critical dimensions and positions to reduce the loss of information caused by interpolation. As illustrated in Fig.\ref{fig:non-ft}, this strategy allows us to increase the LLaMA2 context window from 4k up to 32k without necessitating any fine-tuning. On the other hand, alternative techniques like PI and non-uniform NTK or YaRN result in a significant rise in perplexity beyond a 2$\times$ extension.

\subsection{Extending LLM Context Window to 2048K}
\label{sec:extendto2048k}

\textbf{Progressive extension to 2048k}. Here, we present our approach for expanding the context window of pre-trained LLMs from the standard 4k up to more than 2048k. Our non-uniform positional interpolation, as shown previously, enables an 8$\times$ increase without the need for fine-tuning. However, to achieve greater expansions, such as a 512$\times$ increase, fine-tuning becomes essential. A commonly used strategy involves identifying appropriate RoPE rescaling factors tailored to the desired 2048k context length, followed by subsequent fine-tuning. This process, nonetheless, is hindered by the substantial computational costs required. Furthermore, our findings indicate that successfully fine-tuning LLMs for such large extension ratios poses significant difficulties (refer to Appendix for details).

Fortunately, LongRoPE demonstrates efficacy with both the base model and the fine-tuned version of the extended LLM. Consequently, we present a resource-efficient, incremental approach that reaches the desired 2048k target by performing only 1k fine-tuning iterations, all constrained within a maximum training length of 256k.

$\Diamond$ \textit{LongRoPE search for adapting pre-trained LLM to 256k tokens}. Using LLaMA2 as a case study, we perform a search procedure for intended context window lengths of 128k and 256k tokens. In this process, the model’s receptive field is scaled by factors of 32$\times$ and 64$\times$, correspondingly.

$\Diamond$ \textit{Fine-tuning to 256k}. In this stage, we adapt the pre-trained LLM to support a 256k context window. Initially, LLaMA2 undergoes 400 fine-tuning steps using RoPE rescaling configured for 128k. Afterwards, the rescaling factors are switched to those for 256k at the resultant checkpoint, followed by 600 more fine-tuning steps. This approach offers greater efficiency compared to immediately fine-tuning for the 256k setting.

$\Diamond$ \textit{Scaling fine-tuned extended LLM to 2048k via LongRoPE search}. In the last stage, we apply a subsequent search procedure to the LLM that has already been fine-tuned for 256k-length sequences. This approach increases the context window size dramatically to 2048k, achieving this scale without any additional fine-tuning. Consequently, the overall extension factor reaches $512\times$.

\textbf{Restoration of original context window performance}.  Upon increasing the context window to a substantial 2048k length, a decline in model accuracy is observed when evaluating within the initial context window range. This phenomenon has been documented in the literature on positional interpolation~\cite{pi}; specifically, when high-dimensional position embeddings are compressed into a limited interval corresponding to the original window, the language model's effectiveness deteriorates. Using an extension factor of 512$\times$ causes the position embeddings for the initial 4k window to cluster densely, exacerbating the issue.

To address this issue, we conduct an additional evolutionary search on the expanded LLM, tuning the RoPE rescale parameters specifically for shorter context windows (such as 4k and 8k). Since these lengths necessitate less positional interpolation, we constrain the maximum $\lambda$ considered during the search process. At inference time, the LLM adapts the relevant RoPE rescale parameters dynamically.

 \section{Experiments}

\label{sec:exp}
\subsection{Setup}
\textbf{Evaluation Tasks and Models}. LongRoPE is implemented on both LLaMA2-7B and Mistral-7B architectures, with performance assessed across three key dimensions: (1) perplexity scores when processing lengthy documents, reflecting how well extended-context LLMs handle long inputs; (2) the Passkey retrieval task, designed to test the model’s capability to extract a specific passkey hidden among extensive irrelevant information; and (3) conventional LLM evaluation benchmarks conducted using a context window restricted to 4096 tokens.

\textbf{Fine-tuning}. The LLaMA2 model is fine-tuned using a learning rate of 2e-5 with a linear decay schedule and a global batch size of 32. Training occurs over 400 steps on the Redpajama~\cite{redpajama} dataset, which is segmented into 128k-token chunks and framed with BOS and EOS tokens. Following this, the model is further fine-tuned for an additional 600 steps, starting from the prior checkpoint, to extend to a 256k-token context window. The 128k context configuration utilizes 8 A100 GPUs within a distributed setup~\cite{cube}, whereas scaling to the 256k context requires 16 A100 GPUs. For Mistral, a fixed learning rate of 1e-6 and a global batch size of 64 are applied. Both the 128k and 256k context models adhere to YaRN~\cite{yarn} configurations, involving 400 steps of training on the Together Computer's Long-Data Collections~\cite{mistraldata} with a sequence length of 16k, executed on 4 A100 GPUs.

\textbf{Search}. For target windows not exceeding 256k, the parameters are set as follows: $P$=64, $N_1$=$N_2$=16, $p$=0.3, $\mathcal{T}$=40, with the top 32 individuals chosen for mutation and crossover in every cycle. Perplexity assessments utilize 5 randomly drawn samples from the PG19 validation dataset, each meeting or exceeding the designated context length. When the target window exceeds 512k, both the population size and the sizes for mutation and crossover are reduced by half. In these cases, perplexity is evaluated using 3 randomly selected samples from the Pile-Books3~\cite{pile} validation data.

\textbf{Baselines}. To attain a 2048k context window, we fine-tuned our models starting from context lengths of 128k and 256k. This process produces LongRoPE-2048k (ft=128k) for LLaMA2 and LongRoPE-2048k (ft=256k) for Mistral. We benchmark these four models against leading approaches for expanding context windows, focusing on open-source LLMs that have undergone fine-tuning following positional interpolation techniques such as PI, NTK, and YaRN. The comparative models include Together-32k~\cite{together}, Code LLaMA~\cite{codellama}, LongLoRA-full-FT-100k~\cite{longlora}, YaRN-LLaMA, and YaRN-Mistral~\cite{yarn}.

\subsection{Main Results}

\label{sec:mainresult}
\textbf{Long sequence language modeling within 256k}. Our analysis commences with a comparison to the leading extended LLMs over an evaluation sequence length of 256k tokens. To illustrate the robustness of our approach, we assess performance on two datasets: the Proof-pile~\cite{pg19} and the PG19~\cite{pile} test splits. Perplexity is measured at different context lengths utilizing a sliding window of 256 tokens. For PG19, the evaluation uses the entire set of 100 test documents. For Proof-pile, we adopt the sampling strategy of YaRN~\cite{yarn}, randomly choosing 10 samples, each containing no fewer than 128k tokens.

Table~\ref{tbl:proof-pile} and Table~\ref{tbl:pg19} present a perplexity comparison for LLaMA2 and Mistral models, each extended using various interpolation strategies, on the Proof-pile and PG19 datasets, respectively. Two principal findings emerge: \textbf{(1)} Our extended models consistently demonstrate reduced perplexity as evaluation length increases from 4k to 256k, indicating enhanced capability to utilize extended context. \textbf{(2)} Notably, despite employing a context window that is 16$\times$ longer—an ordinarily difficult scenario for preserving performance at shorter spans—our LongRoPE-2048k models surpass the current leading baselines within a 256k context length.

\textbf{Modeling of Extended Language Sequences Beyond 2000k}. To assess performance on ultra-long textual data, we utilize the Books3~\cite{pile} dataset. For efficient evaluation, a random sample of 20 books—each longer than 2048k tokens—is chosen, applying a 256k-token sliding window approach.

Table~\ref{tbl:books3} demonstrates that LongRoPE is capable of extending the context window of both LLaMA2-7B and Mistral-7B to 2048k tokens, while maintaining perplexity on par with or better than baseline models at shorter context lengths ranging from 8k to 128k. Notably, the performance disparities between the 2048k variants of LLaMA2 and Mistral are significant. While Mistral shows superior results compared to baselines at lower context lengths, its perplexity rises above 7 for contexts longer than 256k. In the case of LLaMA2, performance trends follow expectations, with perplexity generally declining as the context window grows, apart from minor increases observed at the 1024k and 2048k lengths. Furthermore, for LLaMA2, the LongRoPE-2048k model benefits from being fine-tuned at a window size of 256k instead of 128k, a result attributable to the lower secondary extension ratio (8$\times$ compared to 16$\times$). By contrast, Mistral achieves better results with a fine-tuning window of 128k. This difference stems mainly from the application of YaRN's protocol during Mistral's 128k and 256k fine-tuning, where a training length of 16k is used, limiting Mistral's subsequent context window extension capabilities.

\textbf{Passkey retrieval}. Next, we investigate the impact of context window size on generative task performance. To this end, we utilize the passkey retrieval benchmark, as introduced by~\cite{passkey}, which involves extracting a randomly embedded five-digit passkey from a lengthy document. The structure of the prompt employed is outlined in the appendix. For this experiment, we conduct ten repetitions, each time embedding the passkey at a uniformly randomized position within the designated context length.

Figure~\ref{fig:passkey} presents a comparison of retrieval accuracy against baseline models. The accuracy of prior methods declines sharply to zero when surpassing the 128k token mark. Conversely, LongRoPE-LLaMA2-2048k (ft=256k) demonstrates strong performance, consistently achieving high retrieval accuracy ($\ge$90\%) across a wide span, from 4k to 2048k tokens, despite the increased difficulty of extracting a passkey at the million-token scale. Meanwhile, LongRoPE-Mistral-2048k (ft=128k) maintains perfect (100\%) accuracy up to 1800k tokens, with accuracy decreasing to 60\% at 2048k. This behavior matches observations from Table~\ref{tbl:books3}, which shows a moderate rise in perplexity at the 2048k length.

\textbf{Standard benchmarks within original context window}. The LongRoPE-2048k models are assessed on their native context window using the Hugging Face Open LLM Leaderboard~\cite{huggingfaceboard}, with evaluations conducted in both zero-shot and few-shot scenarios. Specifically, the benchmarks include 25-shot ARC-Challenge~\cite{allenai:arc}, 10-shot HellaSwag~\cite{zellers2019hellaswag}, 5-shot MMLU~\cite{mmlu}, and 0-shot TruthfulQA~\cite{lin2021truthfulqa}.

Table~\ref{tbl:commonsense} illustrates that our models perform similarly to prior results on the original benchmark, which targets a shorter context window, and surpass the original Mistral on TruthfulQA by a margin of +0.5\%. Although LongRoPE-LLaMA2-2048k, which was fine-tuned at 256k, experiences somewhat greater performance drop, it generally maintains results within acceptable limits across the evaluated tasks.

\subsection{Ablation Results}

\textbf{Effectiveness of the second positional interpolation}. Within our progressive extension framework, we employ our search algorithm to apply a second round of non-uniform positional interpolation to the fine-tuned extended LLMs. To assess the impact of this approach, we conduct experiments utilizing our fine-tuned LLaMA2-256k model. This model is further extended to sequence lengths of 512k, 1024k, and 2048k through both PI and YaRN methods. As presented in Table~\ref{tbl:secondsearch}, our approach with non-uniform positional interpolation maintains stable perplexity across various extension scales. Conversely, PI and YaRN exhibit rapidly increasing perplexity as the extension ratio grows.

\textbf{Recovery performance at limited context lengths.} To address the drop in performance observed at shorter context windows, we reoptimize the RoPE factors for LongRoPE-2048k using our search algorithm. In this adjustment, we lower the maximum scale factors permitted during the search process, which promotes reduced interpolation for context lengths of 4k and 8k. Table~\ref{tbl:shortperformance} presents a comparison of perplexity metrics for LongRoPE-LLaMA2-2048k on Proof-pile at the 4k and 8k lengths, as well as the mean accuracy across LLM benchmarks. The data reveal a marked improvement in performance for these shorter context windows.

\textbf{Analysis on the two forms of non-uniformities}. To better understand the individual contributions of the two types of non-uniformities, we conduct an ablation study. Specifically, we design two sets of experiments: (i) we extend LLaMA2-7B to shorter sequence lengths of 16k and 32k using various approaches—PI, optimization of the RoPE dimension alone, and optimization of both non-uniformities; (ii) we further extend our previously fine-tuned LLaMA2 with a context length of 256k up to 2048k using the same strategies. We report perplexity without any additional fine-tuning. As demonstrated in Table~\ref{tbl:searchdimension}, introducing non-uniformity into the RoPE dimension yields a substantial perplexity reduction over linear interpolation via PI. Moreover, incorporating non-uniformity in token position leads to better results at 16k and 32k sequence lengths, although this benefit does not extend to the much longer 2048k setting—possibly a consequence of the sheer length involved. At such extreme scales, simply retaining the initial tokens without performing interpolation loses effectiveness, which we highlight as an open direction for future research.

\section{Related Works}

Beyond techniques grounded in position interpolation, this section examines additional related methodologies.

\textbf{Retrieval-based} strategies employ external memory components to store extensive prior context and utilize retrieval mechanisms to obtain relevant documents during inference~\cite{longllama,wang2023augmenting,borgeaud2022improving}. These approaches generally require direct architectural changes to the underlying LLMs. By contrast, our approach is more lightweight, involving only minimal adjustments to the positional embedding scheme. Furthermore, our method extends to a broader set of long-context tasks that go beyond retrieval, including applications such as summarizing lengthy documents and performing few-shot learning.

\textbf{Attention-based context window extensions}. In addition to interpolating positional embeddings, several studies extend the effective context length for inputs by modifying the attention mechanism itself within the inherent context window of the original LLM~\cite{han2023lminfinite, streamingllm,parallelwindow}. The central concept is to address the problem of attention explosion from novel positional tokens through innovative attention masking strategies. These techniques complement positional interpolation approaches.

\textbf{Fine-tuning based approaches} investigate strategies for adapting pre-trained LLMs to longer contexts by adjusting position embeddings during fine-tuning. For instance, models such as Code LLaMA~\cite{codellama}, LLaMA2 Long~\cite{xiong2023effective}, and ScaledRoPE~\cite{liu2023scaling} select an extremely large base value for RoPE and conduct fine-tuning at the intended context length. The approach presented in this work is more adaptable, supporting a range of target lengths and extending to over 2M tokens. Given the significant GPU resource requirements for fine-tuning at extended context lengths (exceeding 128k), more recent approaches like LongLoRA~\cite{longlora} and PoSE~\cite{zhu2023pose} have emerged to alleviate such computational burdens. Our approach operates independently and can be combined with these efficient fine-tuning techniques.

\section{Conclusion}

In this study, we introduce LongRoPE, a technique that significantly expands the context window of LLMs to a remarkable 2048k tokens, while preserving their original performance on shorter sequences. By leveraging two distinct types of non-uniformities present in RoPE positional embeddings and employing an efficient evolutionary search, our approach achieves two primary advantages: it yields optimal initialization for subsequent fine-tuning and permits an 8$\times$ increase in context window size even without any fine-tuning. Further, we design a progressive extension method that utilizes LLMs fine-tuned on 256k-length sequences to attain a context window of 2048k, eliminating the need for additional fine-tuning. Comprehensive experiments demonstrate the strong effectiveness of LongRoPE. We anticipate that the LongRoPE-2048k models will support a wide range of long-context tasks and stimulate future investigation in this area.

\section*{Broader Impacts} The research detailed in this paper is aimed at contributing to the progress of the Machine Learning discipline. While our findings could have a variety of societal implications, we do not identify any that require particular emphasis at this time.

	\balance

\end{document}